\documentclass{report}
\usepackage{amsmath,amssymb}
\setlength{\parindent}{0mm}
\setlength{\parskip}{1em}
\begin{document}
\begin{center}
\rule{6in}{1pt} \
{\large Bo Song \\
{\bf New Coarse Space Components for Additive Schwarz}}

Section of Mathematics \\ University of Geneva \\ 2-4 rue du Lievre \\ CP 64 \\ CH-1211 Geneva \\ Switzerland
\\
{\tt songxjtu@gmail.com}\\
Martin Gander\\
Bo Song\end{center}

The Additive Schwarz method (AS) does not converge in general when used
as a stationary iterative method, and it can only be used as a
preconditioner for a Krylov method. In the two level variant of AS, a
coarse grid correction is added to make the method scalable. We propose
new coarse space components which allow the two level method to become
convergent when used as a stationary iterative method, and show that a
suitable choice makes the method even nil-potent, i.e. it converges in
two iterations, independently of the overlap and the number of
subdomains. One of the coarse spaces we obtain is optimal, which means it
uses the fewest number of basis functions possible for the method to
become nil-potent. We also compare our coarse spaces with GenEO recently
proposed by Spillane, Dolean, Hauret, and et al. We finally illustrate
our theoretical results with numerical experiments.


\end{document}
